% ============ §7 数词 ============
\hdr{§7}{数词：数数、问价、问时间、说第几}

\callout{\textbf{这一节的坐标}：数词 p2--p4（1--1000 与读长数字）、p5--p10（сто́ит／сто́ят 与 юа́нь 的三种形式）、p11--p17（час／мину́та 与时刻）、p18--p22（序数词）、p23--p25（星期与月份）。\\
\textbf{这一节要解决的事}：读出 1949 与 2026；说对「52 元」和「75 元」；说对「2 点 42 分」；把「第 23 路公交车」说成俄语。}

\subhdr{7.1　1--20：先看 13--19 的构造}

数词 p2 的原表就是 1--20。它们的构造有两条规律：

\begin{minipage}{\linewidth}
\tabtop
\begin{tabular}{@{}>{\raggedright\arraybackslash}p{0.08\linewidth}>{\raggedright\arraybackslash}p{0.20\linewidth}>{\raggedright\arraybackslash}p{0.08\linewidth}>{\raggedright\arraybackslash}p{0.20\linewidth}>{\raggedright\arraybackslash}p{0.08\linewidth}>{\raggedright\arraybackslash}p{0.20\linewidth}@{}}
\thead{数} & \thead{俄语} & \thead{数} & \thead{俄语} & \thead{数} & \thead{俄语} \\
\end{tabular}
\tabmid
\begin{tabular}{@{}>{\raggedright\arraybackslash}p{0.08\linewidth}>{\raggedright\arraybackslash}p{0.20\linewidth}>{\raggedright\arraybackslash}p{0.08\linewidth}>{\raggedright\arraybackslash}p{0.20\linewidth}>{\raggedright\arraybackslash}p{0.08\linewidth}>{\raggedright\arraybackslash}p{0.20\linewidth}@{}}
1 & оди́н & 8 & во́семь & 15 & пятна́дцать \\
2 & два & 9 & де́вять & 16 & шестна́дцать \\
3 & три & 10 & де́сять & 17 & семна́дцать \\
4 & четы́ре & 11 & оди́ннадцать & 18 & восемна́дцать \\
5 & пять & 12 & двена́дцать & 19 & девятна́дцать \\
6 & шесть & 13 & трина́дцать & 20 & два́дцать \\
7 & семь & 14 & четы́рнадцать & — & — \\
\end{tabular}
\tabbot
\end{minipage}

\textbf{这张表在说什么}：11--19 全部是「个位数 ＋ \textbf{-надцать}」，只有两处要小心：\textbf{14 четы́рнадцать} 里的 четы́ре 掉了 е（不是 *четы́ренадцать），\textbf{8 восемна́дцать} 与 \textbf{18 восемна́дцать} 里的 во́семь 掉了 ь。

\textbf{逐行读法}：19 是 девятна́дцать（де́вять ＋ надцать，ть 变 т）；其余照拼即可。另外 \textbf{два} 有一个阴性形式 \textbf{две}——数阴性名词时用 две（две ты́сячи、две мину́ты），这条在 7.3、7.4 立刻要用。

\subhdr{7.2　20--1000：整十、整百、整千}

数词 p3 把 20--90、100--900、1000 列在一起。它们的重音大多落在词尾或倒数第二音节，要连着读：

\begin{minipage}{\linewidth}
\tabtop
\begin{tabular}{@{}>{\raggedright\arraybackslash}p{0.22\linewidth}>{\raggedright\arraybackslash}p{0.22\linewidth}>{\raggedright\arraybackslash}p{0.22\linewidth}>{\raggedright\arraybackslash}p{0.26\linewidth}@{}}
\thead{20--90} & \thead{100--900} & \thead{1000} & \thead{读法要点} \\
\end{tabular}
\tabmid
\begin{tabular}{@{}>{\raggedright\arraybackslash}p{0.22\linewidth}>{\raggedright\arraybackslash}p{0.22\linewidth}>{\raggedright\arraybackslash}p{0.22\linewidth}>{\raggedright\arraybackslash}p{0.26\linewidth}@{}}
два́дцать（20） & сто（100） & ты́сяча（1000） & 40 是 со́рок，不规则 \\
три́дцать（30） & две́сти（200） & — & 90 是 девяно́сто \\
пятьдеся́т（50） & три́ста（300） & — & 50--80 是「个位 ＋ де́сят」 \\
шестьдеся́т（60） & четы́реста（400） & — & 500--900 是「个位 ＋ со́т」 \\
се́мьдесят（70） & пятьсо́т（500） & — & ты́сяча 是阴性名词 \\
во́семьдесят（80） & шестьсо́т（600） & — & — \\
девяно́сто（90） & семьсо́т（700） & — & — \\
— & восемьсо́т（800） & — & — \\
— & девятьсо́т（900） & — & — \\
\end{tabular}
\tabbot
\end{minipage}

\textbf{这张表在说什么}：40（со́рок）、90（девяно́сто）、100（сто）三个是不按构造来的，只能记；其余整十、整百都是「个位 ＋ 词根」拼起来的。

\textbf{逐行读法}：50--80 用的是 \textbf{-деся́т}／\textbf{-десят}（пятьдеся́т、шестьдеся́т、се́мьдесят、во́семьдесят），注意重音位置不一样；500--900 用的是 \textbf{-со́т}／\textbf{-сот}（пятьсо́т、шестьсо́т、семьсо́т、восемьсо́т、девятьсо́т）。

\subhdr{7.3　读长数字：从大往小逐段读}

数词 p4 给了四个样本，这是本节最实用的一张：

\begin{itemize}
\item \textbf{123} ＝ сто два́дцать три（100 ＋ 20 ＋ 3）
\item \textbf{564} ＝ пятьсо́т шестьдеся́т четы́ре（500 ＋ 60 ＋ 4）
\item \textbf{1949} ＝ ты́сяча девятьсо́т со́рок де́вять（1000 ＋ 900 ＋ 40 ＋ 9）
\item \textbf{2026} ＝ \textbf{две ты́сячи} два́дцать шесть（2000 ＋ 20 ＋ 6）
\end{itemize}

\textbf{关键一步在第 4 个}：2000 说 \textbf{две ты́сячи}，不是 *два ты́сячи——因为 \textbf{ты́сяча 是阴性名词}，所以「二」用阴性形式 \textbf{две}，而且 2--4 后面的 ты́сяча 要用单数第二格（ты́сячи）。

\aside{这条正是 7.4 那条「2／3／4 ＋ 单数第二格」的同一个规则，只不过这次用在「千」上。所以 7.3 与 7.4 是同一件事，不是两个记忆点。}

\subhdr{7.4　价格：Ско́лько сто́ит? 与 юа́нь 的三种形式}

数词 p5--p6 给了问价句和 юа́нь 的三种形式。这是全篇最能体现「数词决定名词形式」的地方：

\begin{minipage}{\linewidth}
\tabtop
\begin{tabular}{@{}>{\raggedright\arraybackslash}p{0.30\linewidth}>{\raggedright\arraybackslash}p{0.16\linewidth}>{\raggedright\arraybackslash}p{0.44\linewidth}@{}}
\thead{数字结尾} & \thead{用哪种} & \thead{课件例句} \\
\end{tabular}
\tabmid
\begin{tabular}{@{}>{\raggedright\arraybackslash}p{0.30\linewidth}>{\raggedright\arraybackslash}p{0.16\linewidth}>{\raggedright\arraybackslash}p{0.44\linewidth}@{}}
1、21、31、41 …（结尾 1） & юа́нь & Э́то блю́до сто́ит 21 юа́нь. \\
2、3、4、22、23、24 …（结尾 2／3／4） & юа́ня & Э́тот мяч сто́ит 52 юа́ня. \\
5--20、30、40、50 …（结尾 5--9、0、11--14） & юа́ней & Каранда́ш сто́ит 5 юа́ней. \\
\end{tabular}
\tabbot
\end{minipage}

\textbf{这张表在说什么}：名词形式由\textbf{数字的最后一位}决定，不是由数字大小决定：52 用 юа́ня（结尾 2），75 用 юа́ней（结尾 5），21 用 юа́нь（结尾 1），11 用 юа́ней（虽然是 1 结尾，但 11--14 整体归入第三行）。

\textbf{逐行读法}：第一行用\textbf{单数第一格}（юа́нь）；第二行用\textbf{单数第二格}（юа́ня）；第三行用\textbf{复数第二格}（юа́ней）。动词跟着「被数的东西」走：单数用 \textbf{сто́ит}，复数用 \textbf{сто́ят}——\textbf{Моро́женое сто́ит 4 юа́ня.} ／ \textbf{Брю́ки сто́ят 120 юа́ней.}

\textbf{补：为什么 2／3／4 与 5 以上不一样}。这不是随意规定：俄语里 2、3、4 后面跟\textbf{单数第二格}，5 以上跟\textbf{复数第二格}，11--14 因为历史上是「一个十加上几」的复合数，也跟着复数第二格走。所以记忆口诀是三句：\textbf{一用原形，二三四用单二，五以上用复二}。

\warnbox{\textbf{边界}：这条规则只用于「可数名词 ＋ 数词」的搭配。主语若是复数名词（брю́ки、я́блоки），动词用 \textbf{сто́ят}，与数字无关——\textbf{Брю́ки сто́ят 120 юа́ней} 里 сто́ят 是被 брю́ки 决定的。}

\subhdr{7.5　时间：Ско́лько вре́мени? 与 час／мину́та 的三种形式}

数词 p11--p17 给了两个词的三种形式，规则与 юа́нь 一模一样（7.4 的三句口诀照用）：

\begin{minipage}{\linewidth}
\tabtop
\begin{tabular}{@{}>{\raggedright\arraybackslash}p{0.16\linewidth}>{\raggedright\arraybackslash}p{0.20\linewidth}>{\raggedright\arraybackslash}p{0.20\linewidth}>{\raggedright\arraybackslash}p{0.32\linewidth}@{}}
\thead{数字} & \thead{час} & \thead{мину́та} & \thead{例} \\
\end{tabular}
\tabmid
\begin{tabular}{@{}>{\raggedright\arraybackslash}p{0.16\linewidth}>{\raggedright\arraybackslash}p{0.20\linewidth}>{\raggedright\arraybackslash}p{0.20\linewidth}>{\raggedright\arraybackslash}p{0.32\linewidth}@{}}
1 & оди́н час & одна́ мину́та & Сейча́с 11 часо́в 31 мину́та. \\
2、3、4 & два／три／четы́ре часа́ & две／три／четы́ре мину́ты & Сейча́с 2 часа́ 42 мину́ты. \\
5 以上 & пять … часо́в & пять … мину́т & Сейча́с 22 часа́ 15 мину́т. \\
\end{tabular}
\tabbot
\end{minipage}

\textbf{这张表在说什么}：\textbf{час} 的三种形式是 час／часа́／часо́в，\textbf{мину́та} 的是 мину́та／мину́ты／мину́т。三种形式的分工与 юа́нь 完全同构。

\textbf{逐行读法}：第三行的例子最能说明问题——22 点说 \textbf{22 часа́}（结尾 2 → 单数第二格），15 分说 \textbf{15 мину́т}（结尾 5 → 复数第二格）。同一句里两个数词用了两种形式。

\textbf{补：24 小时制与 12 小时制都行}（数词 p11）：\textbf{14 часо́в 45 мину́т} ＝ \textbf{2 часа́ 45 мину́т}；\textbf{21 час} ＝ \textbf{9 часо́в}。问时间三种说法同义：\textbf{Ско́лько вре́мени?} ＝ \textbf{Кото́рый час?} ＝ \textbf{Ско́лько сейча́с вре́мени?}

\subhdr{7.6　序数词、星期与月份}

序数词回答 \textbf{Како́й?}（第几），1--10 是（数词 p18）：

\begin{center}
пе́рвый ｜ второ́й ｜ тре́тий ｜ четвёртый ｜ пя́тый ｜ шесто́й ｜ седьмо́й ｜ восьмо́й ｜ девя́тый ｜ деся́тый
\end{center}

\textbf{怎么用}：第一，复合序数词只有最后一个词变序数（21 ＝ \textbf{два́дцать пе́рвый}；532 ＝ \textbf{пятьсо́т три́дцать второ́й}）；第二，序数词是形容词，四个形式跟名词走（\textbf{девя́тое зада́ние}）；第三，也可以改用「номер ＋ 基数词」：\textbf{зада́ние но́мер де́вять}、\textbf{автобу́с но́мер два́дцать три}（数词 p21--p22）。

\textbf{星期}（数词 p23）：понеде́льник、вто́рник、среда́、четве́рг、пя́тница、суббо́та、воскресе́нье。说「在星期几」用 \textbf{в／во ＋ 第四格}：в понеде́льник、во вто́рник、в сре́ду、в четве́рг、в пя́тницу、в суббо́ту、в воскресе́нье。

\textbf{月份}（数词 p24）：янва́рь、февра́ль、март、апре́ль、май、ию́нь、ию́ль、а́вгуст、сентя́брь、октя́брь、ноя́брь、дека́брь。冬季 дека́брь／янва́рь／февра́ль，春季 март／апре́ль／май，夏季 ию́нь／ию́ль／а́вгуст，秋季 сентя́брь／октя́брь／ноя́брь。

\subhdr{7.7　练习：四个小题}

\begin{enumerate}[label=\textbf{第 \arabic* 题}]
\item （数词 p4）读出来：1949 ｜ 2026。
\item （数词 p6）填 юа́нь／юа́ня／юа́ней：21 \blank；52 \blank；75 \blank；11 \blank。
\item （数词 p13）补词尾：17 час\blank 35 минут\blank；3 час\blank 18 минут\blank。
\item （数词 p22）按模型说「第 25 栋楼」与「第 3 路电车」。
\end{enumerate}

\textbf{第 1 题}　关键一步：从大往小逐段读；2000 里面 ты́сяча 是阴性，用 две。答案：ты́сяча девятьсо́т со́рок де́вять ｜ две ты́сячи два́дцать шесть。

\textbf{第 2 题}　关键一步：看最后一位（11--14 归第三行）。答案：юа́нь ｜ юа́ня ｜ юа́ней ｜ юа́ней。

\textbf{第 3 题}　关键一步：17 结尾 7 → часо́в；35 结尾 5 → мину́т；3 → часа́；18 结尾 8 → мину́т。答案：часо́в ＋ мину́т ｜ часа́ ＋ мину́т。

\textbf{第 4 题}　关键一步：复合序数词只变最后一个词（25 → два́дцать пя́тый）；3 → тре́тий。答案：Э́то два́дцать пя́тый дом.（＝ дом но́мер два́дцать пять）｜ Э́то тре́тий трамва́й.（＝ трамва́й но́мер три）
